\section{Contexto del problema y motivación}

La detección automática de señales de salud mental en redes sociales responde a una necesidad dual. Por una parte, existe una motivación de salud pública: la ideación suicida y la depresión requieren identificación oportuna para facilitar derivación, monitoreo o intervención. Por otra parte, existe una motivación computacional: convertir lenguaje informal, ruidoso y altamente contextual en representaciones útiles para estimar riesgo clínico de forma robusta y reproducible \citep{abdulsalam2024suicidal, montejo2024survey}.

El problema es especialmente complejo por tres razones. Primero, los textos de redes sociales no equivalen a expedientes clínicos: contienen ironía, jerga, referencias implícitas y variación discursiva significativa. Segundo, la clasificación útil para salud mental no siempre es binaria; en muchos casos importa distinguir cambios de estado, niveles de severidad, síntomas específicos o racionales textuales que expliquen la predicción \citep{dechoudhury2016shifts, gaur2019severity}. Tercero, los modelos con mejor desempeño cuantitativo suelen ser también los menos transparentes, lo que introduce un conflicto entre exactitud estadística y confianza clínica.

En Reddit, estas tensiones se vuelven particularmente visibles. Los subreddits de apoyo permiten estudiar desde discursos generales de malestar hasta comunidades explícitamente asociadas con suicidio, depresión, ansiedad o trastornos de personalidad. Ello ha permitido construir tareas distintas: detección de publicaciones suicidas, pronóstico de transición hacia \textit{r/SuicideWatch}, clasificación de severidad, mapeo a categorías DSM-5 y análisis explicativo con modelos de lenguaje \citep{gaur2018dsm5, gaur2019severity, bauer2024llm}.

Para este proyecto, el interés no se centra en cualquier aproximación de PLN, sino en aquellas que aportan evidencia para dos capacidades concretas. La primera es la \textbf{efectividad predictiva}, entendida como la capacidad de discriminar entre clases de riesgo o severidad con métricas sólidas como exactitud, precisión, recall, F1, AUC o variantes ordinales. La segunda es la \textbf{interpretabilidad clínica}, entendida como la posibilidad de relacionar la decisión del sistema con síntomas, categorías diagnósticas, racionales textuales o estructuras conceptuales útiles para especialistas. Esta segunda capacidad es crítica porque el reto académico no busca únicamente maximizar una métrica, sino justificar técnicamente por qué un enfoque podría ser pertinente para un dominio tan sensible.
